% Part: computability
% Chapter: recursive-functions
% Section: partial-functions

\documentclass[../../../include/open-logic-section]{subfiles}

\begin{document}

\olfileid{cmp}{rec}{par}
\olsection{جزوی بازگشتی فنکشن}

بازگشتی فنکشناں دی تعریف دا جواز سمجھّن لئی، غور
کرو کہ قابلِ حساب پر بنیادی بازگشتی نہ ہون والے
فنکشناں دا ساڈا ثبوت اصل وچ ہور وی بہت کجھ ثابت
کردا اے۔ دلیل سادہ سی: اسی صرف ایہہ حقیقت ورتی
کہ فنکشناں $f_0,f_1,\dots$ دی گنتی وار فہرست
ایویں بن سکدی اے کہ $x$ تے $y$ دے فنکشن دے
روپ وچ $f_x(y)$ قابلِ حساب ہووے۔ ایس لئی دلیل
\emph{فنکشناں دی ہر اوہ جماعت اُتے لاگو ہوندی اے
جیہدی فہرست ایس طرح بنائی جا سکدی اے}۔ ایہہ
سانوں مشکل وچ پا دیندا اے: اسی قابلِ حساب فنکشناں
نوں صریح طور اُتے بیان کرنا چاہندے آں؛ پر قابلِ
حساب فنکشناں دے کسے مجموعے دا ایہو جیہا صریح
بیان ساریاں نوں نہیں سمیٹ سکدا!

راہ ایہہ اے کہ \emph{جزوی} فنکشناں نوں وی شامل
کیتا جاوے۔ اسی ویکھاں گے کہ جزوی قابلِ حساب
فنکشناں دی گنتی وار فہرست بناونا \emph{ممکن اے}۔
\iftag{TMs}{ اصل وچ، سانوں کافی حد تک پہلاں ہی
ایہہ پتا اے، کیوں جے ٹیورنگ مشیناں دی فہرست اک
باقاعدہ طریقے نال بنائی جا سکدی اے۔}{} اسی اگے
اپنی قطری دلیل ول واپس آواں گے تے ویکھاں گے کہ
جزوی فنکشن شامل کرن نال اوہ دلیل کیوں نہیں چلدی۔

ہُن سوال ایہہ اے: سارے جزوی بازگشتی فنکشن حاصل
کرن لئی بنیادی بازگشتی فنکشناں وچ کیہ شامل کرنا
پئے گا؟ دو کم کرنے نیں:
\begin{enumerate}
\item بنیادی بازگشتی فنکشناں دی تعریف نوں بدلنا،
تاں کہ جزوی فنکشناں دی وی گنجائش ہووے۔
\item تعریف وچ کجھ \emph{شامل} کرنا، تاں کہ
کجھ نویں جزوی فنکشن وی آ جان۔
\end{enumerate}

پہلا کم سوکھا اے۔ پہلے وانگ، صفر، جانشین تے
پروجیکشن فنکشناں توں شروع کراں گے، تے ترکیب تے
بنیادی بازگشت تحت بندش لواں گے۔ فرق صرف ایہہ
اے کہ ترکیب تے بنیادی بازگشت دی تعریف بدلنی
پئے گی، تاں کہ تعریف وچ آون والے کجھ ٹرم غیر
تعریف شدہ وی ہو سکّن۔ جے $f$ تے $g$ جزوی
فنکشن ہون، تاں $f(x) \fdefined$ دا مطلب ایہہ
لواں گے کہ $f$، $x$ اُتے تعریف شدہ اے، یعنی
$x$، $f$ دے دائرۂ تعریف وچ اے؛ تے $f(x)
\fundefined$ دا اُلٹ مطلب اے، یعنی $f$،~$x$
اُتے تعریف شدہ نہیں۔ $f(x) \simeq g(x)$ دا
مطلب ایہہ لواں گے کہ یا تاں $f(x)$ تے $g(x)$
دونیں غیر تعریف شدہ نیں، یا دونیں تعریف شدہ تے
برابر نیں۔ ہور پیچیدہ ٹرماں لئی وی ایہو علامتاں
ورتاں گے۔ ایہہ طے کر لَیندے آں کہ جے
$h$ تے $g_0$، \dots،~$g_k$ سارے جزوی
فنکشن ہون، تاں
\[
h(g_0(\vec x),\dots,g_k(\vec x))
\]
اُدوں تے صرف اُدوں تعریف شدہ اے جدوں ہر $g_i$،
$\vec x$ اُتے تعریف شدہ ہووے، تے $h$،
$g_0(\vec x)$، \dots،~$g_k(\vec x)$ دلیلاں
دے اِس ٹپل اُتے تعریف شدہ ہووے۔ ایس سمجھ نال جزوی فنکشناں لئی
ترکیب تے بنیادی بازگشت دی تعریف اوپر وانگ ہی اے،
بس «$=$» دی جگہ «$\simeq$» رکھنا پئے گا۔

جزوی فنکشن حاصل کرن لئی بنیادی بازگشتی فنکشناں
دی تعریف وچ جیہڑی شے شامل کراں گے، اوہ
\emph{غیر محدود تلاش دا عامل (unbounded search
  operator)} اے۔ جے $f(x,\vec z)$ قدرتی عدداں
اُتے کوئی وی جزوی فنکشن ہووے، تاں
$\mu x \; f(x,\vec z)$ دی تعریف ایہہ اے:
\begin{quote}
  اوہ کمترین $x$ جیہدے لئی $f(0,\vec z), f(1,\vec z), \dots, f(x,\vec
  z)$ سارے تعریف شدہ ہون، تے $f(x,\vec z) = 0$
  ہووے، جے ایہو جیہا $x$ موجود ہووے
\end{quote}
تے ایہہ سمجھنا اے کہ ہور ہر صورت وچ
$\mu x \; f(x,\vec z)$ غیر تعریف شدہ اے۔ ایس
نال $\mu x \; f(x,\vec z)$ اکّو طرح متعین
ہوندا اے۔

\begin{explain}
غور کرو کہ ساڈی تعریف \iftag{TMs}{ٹیورنگ
مشیناں، یا }{} الگورتھماں یا حساب دے کسے خاص
نمونے دا حوالہ نہیں دیندی۔ پر ترکیب تے بنیادی
بازگشت وانگ، غیر محدود تلاش دے پچھے وی عمل تے
حساب بارے اک بدیہی سمجھ اے۔ جزوی فنکشن دی
حساب پذیری وچ، جیہڑیاں دلیلاں اُتے فنکشن غیر
تعریف شدہ اے، اوہ انہاں انپٹاں دے مطابق نیں
جیہناں لئی حساب نہیں رکدا۔ $\mu x \;
f(x,\vec z)$ حساب کرن دا طریقہ ایہہ ہووے گا:
$f(0,\vec z), f(1,\vec z),
f(2,\vec z)$ حساب کردے جاؤ، جدوں تک صفر
واپس نہ آوے۔ پر جے وچکار دا کوئی حساب نہ رکے،
تاں $\mu x \; f(x,\vec z)$ دا حساب وی نہیں
رکدا۔
\end{explain}

جے $R(x,\vec z)$ کوئی وی تعلق ہووے، تاں
$\mu x \; R(x,\vec z)$ دی تعریف
$\mu x \; (1 \tsub \Char{R}(x,\vec z))$
رکھدے آں۔ دوجے لفظاں وچ، $\mu x \;
R(x,\vec z)$، $x$ دی اوہ کمترین قیمت واپس
دیندا اے جیہدے لئی $R(x,\vec z)$ پورا اترے۔
ایس لئی جے $f(x,\vec z)$ کُل فنکشن ہووے، تاں
$\mu x \; f(x,\vec
z)$، $\mu x \; (f(x,\vec z) = 0)$ دے
برابر اے۔ پر یاد رکھو کہ ساڈی پہلی تعریف ہور
عام اے، کیوں جے اوہ ایہہ گنجائش وی رکھدی اے
کہ $f(x,\vec z)$ ہر تھاں تعریف شدہ نہ ہووے
(ایہدے اُلٹ، تعلق دا خصوصیتی فنکشن ہمیشہ کُل
ہوندا اے)۔

\begin{defn}
\emph{جزوی بازگشتی فنکشناں} دا سیٹ قدرتی عدداں
توں قدرتی عدداں ول مختلف دلیل-گنتیاں والے جزوی
فنکشناں دا سب توں چھوٹا سیٹ اے، جیہدے وچ صفر،
جانشین تے پروجیکشن فنکشن شامل نیں، تے جیہڑا
ترکیب، بنیادی بازگشت تے غیر محدود تلاش تحت بند اے۔
\end{defn}

بے شک، کجھ جزوی بازگشتی فنکشن کُل وی ہون گے،
یعنی ہر دلیل اُتے تعریف شدہ ہون گے۔

\begin{defn}
\ollabel{defn:recursive-fn}
\emph{بازگشتی فنکشناں} دا سیٹ اوہناں جزوی
بازگشتی فنکشناں دا سیٹ اے جیہڑے کُل نیں۔
\end{defn}

بازگشتی فنکشن نوں کدی کدی «کُل بازگشتی»
(total recursive) وی کہندے نیں، تاں کہ اوہدے
ہر تھاں تعریف شدہ ہون اُتے زور دتا جاوے۔

\end{document}
